\section{The stratified nilpotent algebra at projection corank two}
\label{sec:stratified-nilpotent}

The root-free formula of Lemma~\ref{lem:rootfree} contains more
scheme structure than the nilpotency-index statement extracted from it
in the preceding local calculations.  We record that structure before passing
to the exceptional mixed-kernel boundary.

\begin{theorem}[Global residual-colon filtration]
\label{thm:global-residual-colon}
Let \(X=X_R\), let \(\widehat\Delta_R\subset X\) be the reduced
Schubert divisor, and put
\[
 \mathcal J_R^{\mathrm{res}}
 =
 (\mathcal I_{\widehat D_R}:\mathcal I_{\widehat\Delta_R})
 \subset\mathcal O_X.
\]
Then
\begin{equation}\label{eq:global-residual-factorization}
 \mathcal I_{\widehat D_R}
 =
 \mathcal I_{\widehat\Delta_R}\,
 \mathcal J_R^{\mathrm{res}}.
\end{equation}
For every \(j\ge1\),
\begin{equation}\label{eq:global-ann-colon}
 \Ann_{\mathcal O_{\widehat D_R}}
      (\widehat{\mathcal N}_R^{\,j})
 =
 \frac{
   (\mathcal J_R^{\mathrm{res}}:
    \mathcal I_{\widehat\Delta_R}^{\,j-1})
 }{\mathcal I_{\widehat D_R}}.
\end{equation}
Define the closed subscheme of the reduced divisor
\begin{equation}\label{eq:Wj-definition}
 W_j=
 V_{\widehat\Delta_R}\!\left(
  (\mathcal J_R^{\mathrm{res}}:
   \mathcal I_{\widehat\Delta_R}^{\,j-1})
  +\mathcal I_{\widehat\Delta_R}\right).
\end{equation}
If
\[
 \mathcal L=
 \mathcal I_{\widehat\Delta_R}/
 \mathcal I_{\widehat\Delta_R}^{\,2}
\]
is the conormal line of the Schubert divisor, then there are canonical
isomorphisms
\begin{equation}\label{eq:global-graded-piece}
 \widehat{\mathcal N}_R^{\,j}/
 \widehat{\mathcal N}_R^{\,j+1}
 \simeq
 \mathcal L^{\otimes j}\otimes_{\mathcal O_{\widehat\Delta_R}}
 \mathcal O_{W_j}.
\end{equation}
Consequently
\begin{equation}\label{eq:global-gr-colon}
 \gr_{\widehat{\mathcal N}_R}\mathcal O_{\widehat D_R}
 \simeq
 \mathcal O_{\widehat\Delta_R}
 \oplus
 \bigoplus_{j\ge1}
 \mathcal L^{\otimes j}\otimes\mathcal O_{W_j},
\end{equation}
with the multiplication induced by multiplication of powers of
\(\mathcal I_{\widehat\Delta_R}\).  The sum is finite.

These formulas hold on the whole Grassmannian, including every
projection corank and every rank-drop locus.  On a Schur chart with
local equation \(d=\det T\) and residual ideal \(J\), they become
\begin{equation}\label{eq:local-colon-filtration}
 \Ann(N^j)=(J:d^{j-1})/(dJ),\qquad
 N^j/N^{j+1}\simeq
 A/\bigl((d)+(J:d^{j-1})\bigr).
\end{equation}
In particular the nilpotency index at the point is
\begin{equation}\label{eq:nilpotency-colon-index}
 1+\min\{k\ge0:d^k\in J\}.
\end{equation}
\end{theorem}

\begin{proof}
Proposition~\ref{prop:full-fitting} shows locally that
\(\mathcal I_{\widehat D_R}=dJ\), where \(d\) is a local equation of
the Cartier divisor \(\widehat\Delta_R\).  Since \(d\) is a
non-zero-divisor in the regular local ring of \(X\),
\[
 (dJ:d)=J.
\]
Thus \(J\) is the local expression of the intrinsic colon
\(\mathcal J_R^{\mathrm{res}}\), and the local identities glue to
\eqref{eq:global-residual-factorization}.

In \(A/(dJ)\) the nilradical is the image of \((d)\), because the
reduction is \(A/(d)\).  Therefore
\[
 \Ann(N^j)
 =
 (dJ:d^j)/(dJ)
 =
 (J:d^{j-1})/(dJ),
\]
which proves \eqref{eq:global-ann-colon}.  For the graded piece, the
map
\[
 A\longrightarrow N^j/N^{j+1},\qquad
 r\longmapsto r d^j
\]
has kernel
\[
 (d)+(J:d^{j-1}).
\]
Indeed
\[
 rd^j\in d^{j+1}+dJ
 \quad\Longleftrightarrow\quad
 rd^{j-1}\in d^j+J,
\]
and the latter is equivalent to writing
\(r=ds+t\) with \(td^{j-1}\in J\).  This proves the local form of
\eqref{eq:global-graded-piece}.  On overlaps local equations of the
Cartier divisor differ by units, and \(d^j\) transforms by the
\(j\)-th tensor power of the conormal transition function.  Hence the
local identifications glue to \eqref{eq:global-graded-piece} and
\eqref{eq:global-gr-colon}.  Nilpotence of the nilradical makes the
sum finite, and the same cancellation gives
\eqref{eq:nilpotency-colon-index}.
\end{proof}

The theorem converts the globally defined annihilator filtration into
an exact residual-colon filtration.  What remains is to identify the
colon schemes \(W_j\) geometrically on the successive rank strata.
For projection corank one, \(W_1=E_R\) and all higher \(W_j\) are
empty.  The next two results compute \(W_1,W_2\) through the entire
mixed-regular rank-two boundary and then across its first tensor-rank
wall.

Work on a mixed-regular rank-two chart.  Let \(A_0\) be the regular
auxiliary coefficient ring, put
\[
 P=A_0[a,b,c,d],\qquad
 \mathfrak m=(a,b,c,d),\qquad
 \delta=ad-bc,
\]
and localize, or complete, at the point under consideration.  Thus
\[
 \II_{\widehat D_R}=\delta J,\qquad
 J=\delta\mathfrak m+(g_0,\ldots,g_4),
\]
where the \(g_i\) have transverse degree four.  Define the
\emph{unsplit contact ideal}
\begin{equation}\label{eq:contact-saturated-ideal}
 I_H=(\delta,g_0,\ldots,g_4).
\end{equation}
Modulo \(\delta\), the Segre parametrization \(T=uv^t\) identifies
the homogeneous coordinate ring with
\[
 P/(\delta)
 \simeq
 \bigoplus_{n\ge0}
 \Sym^n(\C^2_u)\otimes\Sym^n(\C^2_v),
\]
and the image of \(I_H\) in degree \(n\) is zero for \(n<4\) and
\begin{equation}\label{eq:contact-graded-sections}
 h_H\Sym^{n-4}(\C^2_u)\otimes\Sym^n(\C^2_v)
 \qquad(n\ge4).
\end{equation}

\begin{theorem}[Saturated contact ideal and the fixed vertex component]
\label{thm:saturated-contact}
At every mixed-regular rank-two point one has the exact identities
\begin{equation}\label{eq:J-contact-intersection}
 J=I_H\cap\mathfrak m^3
\end{equation}
and
\begin{equation}\label{eq:failure-contact-vertex}
 \II_{\widehat D_R}
 =\delta I_H\cap\mathfrak m^5.
\end{equation}
If \(h_H\not\equiv0\), the ideal \(I_H\) is
\(\mathfrak m\)-saturated.  Put \(\mathcal C_H=V(I_H)\) for the affine contact cone.  Its
projectivization on the Segre quadric
\[
 Q=\Pj^1_u\times\Pj^1_v\subset\Pj^3
\]
is the effective Cartier divisor
\begin{equation}\label{eq:contact-Cartier-divisor}
 C_H^{\mathrm{proj}}=V(h_H)\times\Pj^1_v
\end{equation}
of class \((4,0)\).  Consequently \(\delta I_H\) is also
\(\mathfrak m\)-saturated, and \(\mathfrak m^5\) in
\eqref{eq:failure-contact-vertex} is an actual irredundant
\(\mathfrak m\)-primary component, not merely a globally defined
ideal.

If
\[
 h_H=\prod_{\nu=1}^s\ell_\nu^{e_\nu},
 \qquad \sum_\nu e_\nu=4,
\]
over an etale coefficient extension, and \(P_\nu\) denotes the
height-two contact prime corresponding to \(\ell_\nu\), then
\begin{equation}\label{eq:associated-primes-mixed-regular}
 \Ass(P/\II_{\widehat D_R})
 =
 \{(\delta),P_1,\ldots,P_s,\mathfrak m\}.
\end{equation}
At the generic point of \(P_\nu\), choose a second regular parameter
\(q_\nu\) so that \(P_\nu=(\delta,q_\nu)\).  The
\(P_\nu\)-primary part of \(\delta I_H\) has the exact two-parameter
normal form
\begin{equation}\label{eq:collision-primary-type}
 Q_{\nu,e_\nu}
 =
 (\delta^2,\delta q_\nu^{e_\nu},q_\nu^{e_\nu+1}).
\end{equation}
For \(e_\nu=1\) this is \(P_\nu^2\); for a collided contact it is the
corresponding non-ordinary primary thickening.

If \(h_H\equiv0\), then
\begin{equation}\label{eq:containment-local-ideal}
 I_H=(\delta),\qquad
 J=\delta\mathfrak m,\qquad
 \II_{\widehat D_R}=\delta^2\mathfrak m
                  =(\delta^2)\cap\mathfrak m^5.
\end{equation}
In particular the rank-two vertex component remains
\(\mathfrak m^5\) on the containment locus.
\end{theorem}

\begin{proof}
The equality \eqref{eq:J-contact-intersection} is a graded
calculation.  The ideal \(I_H\) is generated in transverse degrees
two and four, its only degree-two generator being \(\delta\).
Intersecting with \(\mathfrak m^3\) therefore removes exactly the
scalar multiple of \(\delta\) and leaves
\(\delta\mathfrak m+(g_0,\ldots,g_4)=J\).

Since the initial form of \(\delta\) has degree two and is nonzero in
the polynomial associated graded ring,
\[
 (\mathfrak m^5:\delta)=\mathfrak m^3.
\]
Cancellation in the regular ring gives
\[
 \delta I_H\cap\mathfrak m^5
 =\delta\bigl(I_H\cap(\mathfrak m^5:\delta)\bigr)
 =\delta(I_H\cap\mathfrak m^3)
 =\delta J,
\]
which proves \eqref{eq:failure-contact-vertex}.

Assume first \(h_H\ne0\).  Formula
\eqref{eq:contact-graded-sections} is precisely the section ideal of
\(\OO_Q(-4,0)\subset\OO_Q\), cut out by the nonzero section \(h_H\)
of \(\OO_Q(4,0)\).  Hence it is saturated and defines the Cartier
divisor \eqref{eq:contact-Cartier-divisor}.  The inverse image
\(I_H\subset P\) is therefore \(\mathfrak m\)-saturated.  The product
\(\delta I_H\) is saturated as well: if
\(\mathfrak m^N f\subset\delta I_H\), choose a coordinate not
divisible by the prime element \(\delta\) to see that \(\delta\mid f\);
after cancelling \(\delta\), saturation of \(I_H\) applies.
Thus \(\mathfrak m\notin\Ass(P/\delta I_H)\).
On the other hand \(\delta^2\in\delta I_H\) has transverse degree
four and does not belong to \(\mathfrak m^5\), so the
\(\mathfrak m^5\) factor in
\eqref{eq:failure-contact-vertex} is irredundant.

The divisor \(C_H^{\mathrm{proj}}\) has components
\(e_\nu(\{\ell_\nu=0\}\times\Pj^1)\).  Hence
\(P/I_H\) is Cohen--Macaulay of pure codimension two and its associated
primes are exactly the \(P_\nu\).  Multiplication by \(\delta\) gives
an exact sequence
\[
 0\longrightarrow P/I_H
 \xrightarrow{\ \delta\ } P/\delta I_H
 \longrightarrow P/(\delta)\longrightarrow0.
\]
It follows, and localization shows equality, that the associated
primes of \(P/\delta I_H\) are
\((\delta),P_1,\ldots,P_s\).  Intersecting with the irredundant
\(\mathfrak m\)-primary ideal \(\mathfrak m^5\) gives
\eqref{eq:associated-primes-mixed-regular}.

At the generic point of \(P_\nu\), the Cartier divisor has equation
\(q_\nu^{e_\nu}\), so
\[
 (I_H)_{P_\nu}=(\delta,q_\nu^{e_\nu}),\qquad
 (\delta I_H)_{P_\nu}
 =(\delta^2,\delta q_\nu^{e_\nu}).
\]
The ideal
\[
 (\delta^2,\delta q_\nu^{e_\nu},q_\nu^{e_\nu+1})
\]
is \(P_\nu\)-primary, and its colon by \(\delta\) is
\((\delta,q_\nu^{e_\nu})\).  Therefore
\[
 (\delta)\cap Q_{\nu,e_\nu}
 =
 \delta(\delta,q_\nu^{e_\nu}),
\]
which proves the asserted generic primary normal form.

If \(h_H=0\), every \(g_i\) vanishes modulo \(\delta\).  Since it has
transverse degree four, \(g_i=\delta q_i\) with
\(q_i\in\mathfrak m^2\).  Thus \(I_H=(\delta)\) and
\(J=\delta\mathfrak m\).  Finally
\[
 (\delta^2)\cap\mathfrak m^5
 =\delta^2(\mathfrak m^5:\delta^2)
 =\delta^2\mathfrak m,
\]
again by the polynomial associated graded ring.  This is
\eqref{eq:containment-local-ideal}.
\end{proof}

The preceding theorem gives the requested specialization law without
labelling roots.  A collision changes the Cartier multiplicities in
\(C_H\) and hence the primary type
\eqref{eq:collision-primary-type}; it does not change the vertex
component.  The zero quartic is the unique point at which the Cartier
contact divisor ceases to be a divisor and fills the whole Segre
quadric.

\begin{theorem}[Associated graded nilpotent algebra across coranks one
and two]\label{thm:graded-nilpotent}
On a mixed-regular rank-two chart let
\[
 A=P/(\delta J),\qquad N=\sqrt{0_A}.
\]
Then \(N=(\delta)/(\delta J)\), \(N^3=0\ne N^2\), and
\begin{align}
 N/N^2&\simeq P/I_H, \label{eq:grade-one-contact}\\
 N^2&\simeq P/\mathfrak m. \label{eq:grade-two-vertex}
\end{align}
Equivalently, with \(\mathcal L\) the conormal line of the Schubert
divisor, the associated graded algebra has the intrinsic form
\begin{equation}\label{eq:global-associated-graded}
 \gr_N\OO_{\widehat D_R}
 \simeq
 \OO_{\widehat\Delta_R}
 \oplus
 \bigl(\mathcal L\otimes\OO_{\mathcal C_H}\bigr)
 \oplus
 \bigl(\mathcal L^{\otimes2}\otimes\OO_{\Sigma_R}\bigr)
\end{equation}
on the mixed-regular neighbourhood, with multiplication induced by
the incidence of the affine contact cone \(\mathcal C_H\) with the
rank-two vertex \(\Sigma_R\).
Moreover
\begin{equation}\label{eq:Z1Z2-exact}
 \Ann_A(N)=J/(\delta J),\qquad
 \Ann_A(N^2)=\mathfrak m/(\delta J),
\end{equation}
so
\[
 \ZZ_1=V(I_H\cap\mathfrak m^3),\qquad
 \ZZ_2=V(\mathfrak m).
\]
On the punctured determinant cone, i.e. on the adjacent
projection-corank-one locus, \(\ZZ_1\) is exactly the ramification
support \(E_R\).  Thus \(\mathcal C_H\) is its scheme-theoretic closure away
from the vertex, while \(\ZZ_1\) records in addition the canonical
third-order vertex thickening forced by the Fitting ideal.
\end{theorem}

\begin{proof}
The radical of \(\delta J\) is \((\delta)\), so the nilradical is
generated by the class of \(\delta\).  The isomorphism
\eqref{eq:grade-one-contact} follows by cancelling one factor:
\[
 N/N^2
 =
 \frac{(\delta)}{(\delta^2)+(\delta J)}
 \simeq
 \frac{P}{(\delta)+J}
 =
 \frac{P}{I_H}.
\]
For the square, multiplication by \(\delta^2\) identifies its kernel
with
\[
 (\delta J:\delta^2)=(J:\delta)=\mathfrak m
\]
by Theorem~\ref{thm:collision-depth}; hence
\eqref{eq:grade-two-vertex}.  The transition functions of the local
generator \(\delta\) are those of the conormal line, giving
\eqref{eq:global-associated-graded}.

Cancellation also gives
\[
 (\delta J:\delta)=J,
\]
while the preceding colon calculation gives the second formula in
\eqref{eq:Z1Z2-exact}.  Away from the vertex \(\mathfrak m\) becomes
the unit ideal, so \(J=I_H\).  On the smooth rank-one part of the
determinant divisor, the corank-one normal form
\(t(t,f)\) identifies this contact equation with the ramification
equation.  Hence the restriction is precisely \(E_R\).
\end{proof}

We finally cross the first non-mixed-regular wall.  This is the
rank-one tensor wall in the one-dimensional mixed kernel, explicitly
requested by the referee.

\begin{definition}\label{def:decomposable-kernel-regular}
Assume \(A_H\) is injective and \(\overline B_H\) is surjective with
one-dimensional kernel generated by a decomposable tensor
\(h_1\otimes w_1\).  Choose \(h_2,w_2\) completing bases and put
\[
 q_1=\overline B_H(h_1\otimes w_2),\qquad
 q_2=\overline B_H(h_2\otimes w_1).
\]
We call the point \emph{transverse decomposable-kernel regular} if
\[
 \det[q_1,q_2,\overline C_H(w_1^2)]\ne0.
\]
This condition is independent of the chosen completing vectors.  It
says that the contact root forced by the decomposable kernel is simple
in the direction transverse to the kernel wall.
\end{definition}

\begin{theorem}[The decomposable mixed-kernel wall]
\label{thm:decomposable-kernel-wall}
At a transverse decomposable-kernel regular rank-two point there are
bases for which
\begin{equation}\label{eq:decomposable-linear-block}
 L_{\mathrm{dec}}(T)=
 \begin{pmatrix}
 c&d&0&0\\
 0&0&a&b\\
 0&0&c&d
 \end{pmatrix}.
\end{equation}
Let
\[
 J_{\mathrm{dec}}
 =I_3[L_{\mathrm{dec}}(T),C\,T^{(2)}],\qquad
 K_{\mathrm{dec}}=(c,d,a^2,ab,b^2).
\]
Then
\begin{equation}\label{eq:decomposable-colon}
 (J_{\mathrm{dec}}:\delta)=K_{\mathrm{dec}}.
\end{equation}
Hence the local failure ring still has nilradical of exact index three,
but its second intrinsic depth support is no longer reduced:
\begin{equation}\label{eq:decomposable-Z2}
 \ZZ_2
 =
 V(K_{\mathrm{dec}}).
\end{equation}
Its support is the rank-two vertex, and its transverse local length is
three.  Thus crossing the tensor-rank-two wall changes the
scheme structure of \(\ZZ_2\) before it changes its set-theoretic
support.
\end{theorem}

\begin{proof}
The stated normal form is obtained by taking
\(\ker\overline B_H=\C(h_1\otimes w_1)\) and using
\[
 q_1,\ q_2,\ q_3=\overline B_H(h_2\otimes w_2)
\]
as target basis.  Applying \(1_H\otimes T\) gives
\eqref{eq:decomposable-linear-block}.  Its maximal minors generate
\[
 I_3(L_{\mathrm{dec}})=\delta(c,d).
\]
After replacing the first quadratic target column by its component
modulo \(\langle q_1,q_2\rangle\), the transversality assumption
allows us to normalize
\(\overline C_H(w_1^2)\equiv q_3\).  The minors containing one
column from each mixed group and this quadratic column give
\[
 \delta a^2,\qquad \delta ab,\qquad \delta b^2
 \in J_{\mathrm{dec}}.
\]
Together with \(I_3(L_{\mathrm{dec}})\) this proves
\(\delta K_{\mathrm{dec}}\subset J_{\mathrm{dec}}\), hence
\(K_{\mathrm{dec}}\subset(J_{\mathrm{dec}}:\delta)\).

For the reverse inclusion use the transverse grading.
The degree-three part of \(J_{\mathrm{dec}}\) is exactly
\(\delta(c,d)\); every minor using a quadratic column has degree at
least four.  Modulo \(K_{\mathrm{dec}}\), every class is represented
by
\[
 \lambda+\mu a+\nu b.
\]
If \(r\delta\in J_{\mathrm{dec}}\), the degree-two term rules out
\(\lambda\ne0\), while the degree-three identity
\[
 (\mu a+\nu b)\delta\in\delta(c,d)
\]
forces \(\mu=\nu=0\).  Thus \(r\in K_{\mathrm{dec}}\), proving
\eqref{eq:decomposable-colon}.

Since \(\delta\in(c,d)\), one has
\(\delta^2\in\delta(c,d)\subset J_{\mathrm{dec}}\), whereas
\(\delta\notin J_{\mathrm{dec}}\) by degree.  Therefore the
nilradical of \(P/(\delta J_{\mathrm{dec}})\) has exact index three.
Finally
\[
 \Ann(N^2)
 =(\delta J_{\mathrm{dec}}:\delta^2)
 =(J_{\mathrm{dec}}:\delta)
 =K_{\mathrm{dec}}.
\]
The quotient by \(K_{\mathrm{dec}}\) has basis \(1,a,b\), proving the
length assertion.
\end{proof}

\begin{proposition}[The containment locus is the finite line scheme]
\label{prop:containment-lines}
For fixed \(R\in G_4^\circ\), the locus
\[
 \mathcal C_R=\{H\in\Gr(2,V):h_H\equiv0\}
\]
is canonically the Fano scheme \(F_1(Y_R)\) of lines on the smooth
quartic \(Y_R\).  It is finite and reduced.  At its mixed-regular
points the local failure ideal is
\eqref{eq:containment-local-ideal}, and the contact morphism
\(\RR_R\to\Gr(2,V)\) has fibre \(\Pj^1\), rather than a finite
length-four fibre.
\end{proposition}

\begin{proof}
The line of hyperplanes containing \(H\) is
\(\Pj(H^\perp)\subset\Pj(V^*)\).  Its quartic restriction is
\(h_H\), so \(h_H=0\) identically exactly when
\(\Pj(H^\perp)\subset Y_R\).  This identifies \(\mathcal C_R\) with
\(F_1(Y_R)\).

If \(L\subset Y_R\) is such a line, the normal sequence in
\(\Pj^3\) gives
\[
 \deg N_{L/Y_R}
 =
 \deg N_{L/\Pj^3}-\deg N_{Y_R/\Pj^3}|_L
 =2-4=-2.
\]
Since \(N_{L/Y_R}\) has rank one,
\(N_{L/Y_R}\simeq\OO_{\Pj^1}(-2)\), and therefore
\(H^0(N_{L/Y_R})=0\).  The Zariski tangent space of the Fano scheme
vanishes at every point.  The Fano scheme is projective and
zero-dimensional; a zero-dimensional local algebra with zero tangent
space is a field by Nakayama.  Hence it is finite reduced.  The final
two assertions follow from Theorem~\ref{thm:saturated-contact} and
the definition of the incidence fibre.
\end{proof}
